% --- Archon physics formalization source begin ---
% archon:physics
% archon:covers PhyXMiniProblems/problem_phyx_mini_0594.lean
% archon:source-report reports/phyx_mini/problem_phyx_mini_0594.source.json

\chapter{Physics problem phyx\_mini\_0594}
\label{ch:PhyXMiniProblems_problem_phyx_mini_0594}

\paragraph{Problem source.}
\#\# Physical scenario

Figure shows the path taken by a knot in a jet of ionized gas that has been expelled from a galaxy. The knot travels at constant velocity \$\textbackslash{}vec\{v\}\$ at angle \$\textbackslash{}theta\$ from the direction of Earth. The knot occasionally emits a burst of light, which is eventually detected on Earth. Two bursts are indicated in figure, separated by time \$t\$ as measured in a stationary frame near the bursts.  The apparent distance \$D\_\{app\}\$ traveled by the knot between the two bursts is the distance across an Earth-observer's view of the knot's path. The apparent time \$T\_\{app\}\$ between the bursts is the difference in the arrival times of the light from them. The apparent speed of the knot is then \$V\_\{app\} = D\_\{app\} / T\_\{app\}\$.  When superluminal (faster than light) jets were first observed, they seemed to defy special relativity—at least until the correct geometry was understood.

\#\# Question

Evaluate \$V\_\{app\}\$ for \$v = 0.980c\$ and \$\textbackslash{}theta = 30.0\textasciicircum{}\textbackslash{}circ.\$

\#\# Answer choices

- A: A: 2.75c

- B: B: 1.96c

- C: C: 3.24c

- D: D: 4.50c

\#\# Figure caption (auxiliary; use the image as primary evidence)

The image is a diagram showing the path of a knot of ionized gas in space. Here's a detailed description:

1. **Path Line**: A diagonal straight line represents the path of the ionized gas, marked as "Path of knot of ionized gas."

2. **Bursts**: 
   - Two points labeled as "Burst 1" and "Burst 2" are on the path. They are depicted with red starburst symbols.
   
3. **Velocity Vector**: 
   - A magenta arrow labeled with \textbackslash{}( \textbackslash{}vec\{v\} \textbackslash{}) represents the velocity of the gas moving from Burst 1 to Burst 2.
   
4. **Angle**: 
   - An angle \textbackslash{}( \textbackslash{}theta \textbackslash{}) is formed between the path line and a vertical line extending downward from Burst 1.

5. **Light Rays**:
   - Red arrows depict light rays headed to Earth, originating from each burst.

6. **Apparent Distance**:
   - A horizontal double-headed arrow between the vertical light rays extending from Burst 1 and Burst 2, labeled as \textbackslash{}( D\_\{\textbackslash{}text\{app\}\} \textbackslash{}), indicates the apparent distance.

7. **Text**:
   - "Light rays headed to Earth" is written next to the arrows representing the light rays.

This diagram seems to illustrate the motion and light travel path of a high-speed astrophysical jet, potentially explaining apparent size or position measurements based on light arrival.

\#\# Dataset metadata

Relativity; Physical Model Grounding Reasoning, Spatial Relation Reasoning

\paragraph{Recorded answer/context.}
C: C: 3.24c

\paragraph{Figure/image path.}
/root/proposal\_for\_physic/science-mango/phyx\_mini\_run/phyx\_data/test\_image/594.png

\paragraph{Formalization target.}
create a compiling Lean file with sorry bodies at `PhyXMiniProblems/problem\_phyx\_mini\_0594.lean`. The Lean declarations must preserve the physical quantities, dimensions or dimensional roles, figure labels, governing-law hypotheses, and final relation expressed by this problem.
Use Mathlib/Physlib names found through LeanExplore where available. If a physics API is missing, introduce faithful local abstractions rather than scalar placeholder aliases.

\begin{theorem}[Physics formalization target]
\label{thm:physics:phyx_mini_0594:target}
\lean{PhyXMiniProblems.ProblemPhyXMini0594.apparent_speed_of_ionized_gas_knot_matches_answer_C}
\uses{def:physics:phyx-mini-0594:phyxminiproblems-problemphyxmini0594-speedfractionoflight, def:physics:phyx-mini-0594:phyxminiproblems-problemphyxmini0594-apparentjetmotionsetup, def:physics:phyx-mini-0594:phyxminiproblems-problemphyxmini0594-matchesjetscenario, def:physics:phyx-mini-0594:phyxminiproblems-problemphyxmini0594-matchesproblemreadouts, def:physics:phyx-mini-0594:phyxminiproblems-problemphyxmini0594-matchessuppliedjetfigure, def:physics:phyx-mini-0594:phyxminiproblems-problemphyxmini0594-hasphysicaljetobservationparameters, def:physics:phyx-mini-0594:phyxminiproblems-problemphyxmini0594-satisfiesapparentmotionlaws, def:physics:phyx-mini-0594:phyxminiproblems-problemphyxmini0594-matchesanswertonearesthundredth}
The assigned autoformalize agent should translate this physics problem into Lean declarations in the covered file, with theorem and lemma proof bodies written as `by sorry`.
\end{theorem}
\begin{proof}
This is an autoformalization task, not a proof task. Produce statements that can later be proved by the physics prover without weakening the physical model.
\end{proof}

% --- Archon declaration topology begin ---
\section{Declaration topology}
\begin{definition}[Length Quantity]
  \label{def:physics:phyx-mini-0594:phyxminiproblems-problemphyxmini0594-lengthquantity}
  \lean{PhyXMiniProblems.ProblemPhyXMini0594.LengthQuantity}
  A nonnegative, unit-independent physical length
\end{definition}
\begin{proof}
  This is the declared model definition of Length Quantity.
\end{proof}
\begin{definition}[Duration Quantity]
  \label{def:physics:phyx-mini-0594:phyxminiproblems-problemphyxmini0594-durationquantity}
  \lean{PhyXMiniProblems.ProblemPhyXMini0594.DurationQuantity}
  A nonnegative, unit-independent physical duration
\end{definition}
\begin{proof}
  This is the declared model definition of Duration Quantity.
\end{proof}
\begin{definition}[Speed Quantity]
  \label{def:physics:phyx-mini-0594:phyxminiproblems-problemphyxmini0594-speedquantity}
  \lean{PhyXMiniProblems.ProblemPhyXMini0594.SpeedQuantity}
  A nonnegative physical speed
\end{definition}
\begin{proof}
  This is the declared model definition of Speed Quantity.
\end{proof}
\begin{definition}[length In Meters]
  \label{def:physics:phyx-mini-0594:phyxminiproblems-problemphyxmini0594-lengthinmeters}
  \lean{PhyXMiniProblems.ProblemPhyXMini0594.lengthInMeters}
  \uses{def:physics:phyx-mini-0594:phyxminiproblems-problemphyxmini0594-lengthquantity}
  Metre readout of a physical length
\end{definition}
\begin{proof}
  This is the declared model definition of length In Meters.
\end{proof}
\begin{definition}[duration In Seconds]
  \label{def:physics:phyx-mini-0594:phyxminiproblems-problemphyxmini0594-durationinseconds}
  \lean{PhyXMiniProblems.ProblemPhyXMini0594.durationInSeconds}
  \uses{def:physics:phyx-mini-0594:phyxminiproblems-problemphyxmini0594-durationquantity}
  Second readout of a duration measured in the stationary burst frame
\end{definition}
\begin{proof}
  This is the declared model definition of duration In Seconds.
\end{proof}
\begin{definition}[speed In Meters Per Second]
  \label{def:physics:phyx-mini-0594:phyxminiproblems-problemphyxmini0594-speedinmeterspersecond}
  \lean{PhyXMiniProblems.ProblemPhyXMini0594.speedInMetersPerSecond}
  \uses{def:physics:phyx-mini-0594:phyxminiproblems-problemphyxmini0594-speedquantity}
  Metres-per-second readout of a physical speed
\end{definition}
\begin{proof}
  This is the declared model definition of speed In Meters Per Second.
\end{proof}
\begin{definition}[vacuum Speed Of Light In Meters Per Second]
  \label{def:physics:phyx-mini-0594:phyxminiproblems-problemphyxmini0594-vacuumspeedoflightinmeterspersecond}
  \lean{PhyXMiniProblems.ProblemPhyXMini0594.vacuumSpeedOfLightInMetersPerSecond}
  Physlib's exact vacuum speed of light, read in metres per second
\end{definition}
\begin{proof}
  This is the declared model definition of vacuum Speed Of Light In Meters Per Second.
\end{proof}
\begin{definition}[speed Fraction Of Light]
  \label{def:physics:phyx-mini-0594:phyxminiproblems-problemphyxmini0594-speedfractionoflight}
  \lean{PhyXMiniProblems.ProblemPhyXMini0594.speedFractionOfLight}
  \uses{def:physics:phyx-mini-0594:phyxminiproblems-problemphyxmini0594-speedquantity, def:physics:phyx-mini-0594:phyxminiproblems-problemphyxmini0594-speedinmeterspersecond, def:physics:phyx-mini-0594:phyxminiproblems-problemphyxmini0594-vacuumspeedoflightinmeterspersecond}
  Dimensionless speed as a multiple of the vacuum speed of light
\end{definition}
\begin{proof}
  This is the declared model definition of speed Fraction Of Light.
\end{proof}
\begin{definition}[angle From Degrees]
  \label{def:physics:phyx-mini-0594:phyxminiproblems-problemphyxmini0594-anglefromdegrees}
  \lean{PhyXMiniProblems.ProblemPhyXMini0594.angleFromDegrees}
  Convert a degree readout to Mathlib's type of angles
\end{definition}
\begin{proof}
  This is the declared model definition of angle From Degrees.
\end{proof}
\begin{definition}[Burst Label]
  \label{def:physics:phyx-mini-0594:phyxminiproblems-problemphyxmini0594-burstlabel}
  \lean{PhyXMiniProblems.ProblemPhyXMini0594.BurstLabel}
  The two light-emission events marked by red starbursts in the figure
\end{definition}
\begin{proof}
  This is the declared model definition of Burst Label.
\end{proof}
\begin{definition}[Astrophysical Source]
  \label{def:physics:phyx-mini-0594:phyxminiproblems-problemphyxmini0594-astrophysicalsource}
  \lean{PhyXMiniProblems.ProblemPhyXMini0594.AstrophysicalSource}
  The astrophysical source from which the jet is expelled
\end{definition}
\begin{proof}
  This is the declared model definition of Astrophysical Source.
\end{proof}
\begin{definition}[Moving Object]
  \label{def:physics:phyx-mini-0594:phyxminiproblems-problemphyxmini0594-movingobject}
  \lean{PhyXMiniProblems.ProblemPhyXMini0594.MovingObject}
  The moving object whose apparent speed is observed
\end{definition}
\begin{proof}
  This is the declared model definition of Moving Object.
\end{proof}
\begin{definition}[Knot Motion Model]
  \label{def:physics:phyx-mini-0594:phyxminiproblems-problemphyxmini0594-knotmotionmodel}
  \lean{PhyXMiniProblems.ProblemPhyXMini0594.KnotMotionModel}
  The motion model explicitly stated in the problem
\end{definition}
\begin{proof}
  This is the declared model definition of Knot Motion Model.
\end{proof}
\begin{definition}[Figure Text Label]
  \label{def:physics:phyx-mini-0594:phyxminiproblems-problemphyxmini0594-figuretextlabel}
  \lean{PhyXMiniProblems.ProblemPhyXMini0594.FigureTextLabel}
  Literal mathematical or descriptive labels visible in the primary image
\end{definition}
\begin{proof}
  This is the declared model definition of Figure Text Label.
\end{proof}
\begin{definition}[Figure Arrow]
  \label{def:physics:phyx-mini-0594:phyxminiproblems-problemphyxmini0594-figurearrow}
  \lean{PhyXMiniProblems.ProblemPhyXMini0594.FigureArrow}
  The four salient arrows drawn in the primary image
\end{definition}
\begin{proof}
  This is the declared model definition of Figure Arrow.
\end{proof}
\begin{definition}[Figure Direction]
  \label{def:physics:phyx-mini-0594:phyxminiproblems-problemphyxmini0594-figuredirection}
  \lean{PhyXMiniProblems.ProblemPhyXMini0594.FigureDirection}
  Qualitative directions of arrows relative to the observer and jet path
\end{definition}
\begin{proof}
  This is the declared model definition of Figure Direction.
\end{proof}
\begin{definition}[Superluminal Jet Figure]
  \label{def:physics:phyx-mini-0594:phyxminiproblems-problemphyxmini0594-superluminaljetfigure}
  \lean{PhyXMiniProblems.ProblemPhyXMini0594.SuperluminalJetFigure}
  \uses{def:physics:phyx-mini-0594:phyxminiproblems-problemphyxmini0594-burstlabel, def:physics:phyx-mini-0594:phyxminiproblems-problemphyxmini0594-figuretextlabel, def:physics:phyx-mini-0594:phyxminiproblems-problemphyxmini0594-figurearrow, def:physics:phyx-mini-0594:phyxminiproblems-problemphyxmini0594-figuredirection}
  Semantic transcription of the supplied diagram. These fields record only which objects, labels, and directional relations are drawn; no numerical apparent speed is stored here
\end{definition}
\begin{proof}
  This is the declared model definition of Superluminal Jet Figure.
\end{proof}
\begin{definition}[Apparent Jet Motion Setup]
  \label{def:physics:phyx-mini-0594:phyxminiproblems-problemphyxmini0594-apparentjetmotionsetup}
  \lean{PhyXMiniProblems.ProblemPhyXMini0594.ApparentJetMotionSetup}
  \uses{def:physics:phyx-mini-0594:phyxminiproblems-problemphyxmini0594-lengthquantity, def:physics:phyx-mini-0594:phyxminiproblems-problemphyxmini0594-durationquantity, def:physics:phyx-mini-0594:phyxminiproblems-problemphyxmini0594-speedquantity, def:physics:phyx-mini-0594:phyxminiproblems-problemphyxmini0594-burstlabel, def:physics:phyx-mini-0594:phyxminiproblems-problemphyxmini0594-astrophysicalsource, def:physics:phyx-mini-0594:phyxminiproblems-problemphyxmini0594-movingobject, def:physics:phyx-mini-0594:phyxminiproblems-problemphyxmini0594-knotmotionmodel, def:physics:phyx-mini-0594:phyxminiproblems-problemphyxmini0594-superluminaljetfigure}
  The independent physical quantities in the observation. In particular, apparentSpeed is an observable field; it is not defined from an answer choice or from the requested numerical result
\end{definition}
\begin{proof}
  This is the declared model definition of Apparent Jet Motion Setup.
\end{proof}
\begin{definition}[Matches Jet Scenario]
  \label{def:physics:phyx-mini-0594:phyxminiproblems-problemphyxmini0594-matchesjetscenario}
  \lean{PhyXMiniProblems.ProblemPhyXMini0594.MatchesJetScenario}
  \uses{def:physics:phyx-mini-0594:phyxminiproblems-problemphyxmini0594-apparentjetmotionsetup}
  Categorical physical facts stated in the scenario
\end{definition}
\begin{proof}
  This is the declared model definition of Matches Jet Scenario.
\end{proof}
\begin{definition}[Matches Problem Readouts]
  \label{def:physics:phyx-mini-0594:phyxminiproblems-problemphyxmini0594-matchesproblemreadouts}
  \lean{PhyXMiniProblems.ProblemPhyXMini0594.MatchesProblemReadouts}
  \uses{def:physics:phyx-mini-0594:phyxminiproblems-problemphyxmini0594-speedinmeterspersecond, def:physics:phyx-mini-0594:phyxminiproblems-problemphyxmini0594-vacuumspeedoflightinmeterspersecond, def:physics:phyx-mini-0594:phyxminiproblems-problemphyxmini0594-anglefromdegrees, def:physics:phyx-mini-0594:phyxminiproblems-problemphyxmini0594-apparentjetmotionsetup}
  The two numerical readouts supplied by the problem. No apparent-distance, arrival-time, apparent-speed, or answer-choice value occurs here
\end{definition}
\begin{proof}
  This is the declared model definition of Matches Problem Readouts.
\end{proof}
\begin{definition}[Matches Supplied Jet Figure]
  \label{def:physics:phyx-mini-0594:phyxminiproblems-problemphyxmini0594-matchessuppliedjetfigure}
  \lean{PhyXMiniProblems.ProblemPhyXMini0594.MatchesSuppliedJetFigure}
  \uses{def:physics:phyx-mini-0594:phyxminiproblems-problemphyxmini0594-apparentjetmotionsetup}
  Qualitative geometry read directly from the primary raster
\end{definition}
\begin{proof}
  This is the declared model definition of Matches Supplied Jet Figure.
\end{proof}
\begin{definition}[Has Physical Jet Observation Parameters]
  \label{def:physics:phyx-mini-0594:phyxminiproblems-problemphyxmini0594-hasphysicaljetobservationparameters}
  \lean{PhyXMiniProblems.ProblemPhyXMini0594.HasPhysicalJetObservationParameters}
  \uses{def:physics:phyx-mini-0594:phyxminiproblems-problemphyxmini0594-lengthinmeters, def:physics:phyx-mini-0594:phyxminiproblems-problemphyxmini0594-durationinseconds, def:physics:phyx-mini-0594:phyxminiproblems-problemphyxmini0594-speedinmeterspersecond, def:physics:phyx-mini-0594:phyxminiproblems-problemphyxmini0594-vacuumspeedoflightinmeterspersecond, def:physics:phyx-mini-0594:phyxminiproblems-problemphyxmini0594-apparentjetmotionsetup}
  Physical-domain conditions for the approaching-jet branch depicted in the figure. They state only positivity, acuteness, and subluminal true motion
\end{definition}
\begin{proof}
  This is the declared model definition of Has Physical Jet Observation Parameters.
\end{proof}
\begin{definition}[Satisfies Apparent Motion Laws]
  \label{def:physics:phyx-mini-0594:phyxminiproblems-problemphyxmini0594-satisfiesapparentmotionlaws}
  \lean{PhyXMiniProblems.ProblemPhyXMini0594.SatisfiesApparentMotionLaws}
  \uses{def:physics:phyx-mini-0594:phyxminiproblems-problemphyxmini0594-lengthinmeters, def:physics:phyx-mini-0594:phyxminiproblems-problemphyxmini0594-durationinseconds, def:physics:phyx-mini-0594:phyxminiproblems-problemphyxmini0594-speedinmeterspersecond, def:physics:phyx-mini-0594:phyxminiproblems-problemphyxmini0594-vacuumspeedoflightinmeterspersecond, def:physics:phyx-mini-0594:phyxminiproblems-problemphyxmini0594-apparentjetmotionsetup}
  The governing constant-velocity, projection, and light-arrival relations. The second burst occurs closer to Earth by the earthward projection of the knot's displacement. Its light therefore has less distance to travel, so the arrival interval is the emission interval minus that projection divided by c. These laws contain no numerical value for the requested apparent speed
\end{definition}
\begin{proof}
  This is the declared model definition of Satisfies Apparent Motion Laws.
\end{proof}
\begin{definition}[Answer Choice]
  \label{def:physics:phyx-mini-0594:phyxminiproblems-problemphyxmini0594-answerchoice}
  \lean{PhyXMiniProblems.ProblemPhyXMini0594.AnswerChoice}
  Labels of the four multiple-choice answers
\end{definition}
\begin{proof}
  This is the declared model definition of Answer Choice.
\end{proof}
\begin{definition}[speed Of Light Multiple]
  \label{def:physics:phyx-mini-0594:phyxminiproblems-problemphyxmini0594-answerchoice-speedoflightmultiple}
  \lean{PhyXMiniProblems.ProblemPhyXMini0594.AnswerChoice.speedOfLightMultiple}
  \uses{def:physics:phyx-mini-0594:phyxminiproblems-problemphyxmini0594-answerchoice}
  Multiple of c printed beside each answer label
\end{definition}
\begin{proof}
  This is the declared model definition of speed Of Light Multiple.
\end{proof}
\begin{definition}[Matches Answer To Nearest Hundredth]
  \label{def:physics:phyx-mini-0594:phyxminiproblems-problemphyxmini0594-matchesanswertonearesthundredth}
  \lean{PhyXMiniProblems.ProblemPhyXMini0594.MatchesAnswerToNearestHundredth}
  \uses{def:physics:phyx-mini-0594:phyxminiproblems-problemphyxmini0594-speedquantity, def:physics:phyx-mini-0594:phyxminiproblems-problemphyxmini0594-speedfractionoflight, def:physics:phyx-mini-0594:phyxminiproblems-problemphyxmini0594-answerchoice}
  A speed matches a displayed answer to the nearest hundredth of c when its dimensionless ratio to c differs by at most 0.005
\end{definition}
\begin{proof}
  This is the declared model definition of Matches Answer To Nearest Hundredth.
\end{proof}
\begin{definition}[recorded Dataset Answer]
  \label{def:physics:phyx-mini-0594:phyxminiproblems-problemphyxmini0594-recordeddatasetanswer}
  \lean{PhyXMiniProblems.ProblemPhyXMini0594.recordedDatasetAnswer}
  \uses{def:physics:phyx-mini-0594:phyxminiproblems-problemphyxmini0594-answerchoice}
  The answer label recorded in the supplied dataset
\end{definition}
\begin{proof}
  This is the declared model definition of recorded Dataset Answer.
\end{proof}
% --- Archon declaration topology end ---
% --- Archon physics formalization source end ---
